\documentclass{article}

\usepackage{neurips_2025}

\usepackage[utf8]{inputenc}
\usepackage[T1]{fontenc}
\usepackage{hyperref}
\usepackage{url}
\usepackage{booktabs}
\usepackage{amsfonts}
\usepackage{nicefrac}
\usepackage{microtype}
\usepackage{xcolor}
\usepackage{graphicx}
\usepackage{subfigure}
\usepackage{soul}
\usepackage{enumitem}
\usepackage{adjustbox}
\usepackage{amsmath}
\usepackage{amssymb}

% --- Author comment macros ---
\newcommand{\todo}[1]{\textcolor{red}{[TODO: #1]}}
\newcommand{\andrea}[1]{\textcolor{purple}{[AY: #1]}}
\newcommand{\louise}[1]{\textcolor{blue}{[LY: #1]}}

\title{De-risking Frontier Projects: A Systematic Framework \& Review for Tackling High-Uncertainty Work}

\author{
  Alycia Lee\textsuperscript{1,2,*} \And
  Brando Miranda\textsuperscript{1,2,*} \And
  Sudharsan Sundar\textsuperscript{1} \And
  Sanmi Koyejo\textsuperscript{1} \\
  \textsuperscript{1}Department of Computer Science, Stanford University, Palo Alto, USA \\
  \textsuperscript{*}Corresponding authors: \texttt{alylee15@stanford.edu}, \texttt{brando9@stanford.edu}
}

\begin{document}

\maketitle

\begin{abstract}
The prevailing belief in computer-science training holds that research mastery emerges primarily through prolonged apprenticeship---a model that leaks time, talent, and R\&D (Research and Development) capital.
We challenge this assumption by formalising \textbf{Velocity--Vectoring}, a decision loop distilled from Bernstein's 197 lectures~\cite{bernstein2025cs197}.
The loop ranks uncertainties by potential impact, then resolves the highest-risk item through the leanest viable experiment.
We instantiate the loop as an open-source toolkit comprising assumption-mapping worksheets, one-week sprint planners, and automated ``velocity gauges'' that visualise learning per unit time.
We evaluate the toolkit on three high-uncertainty initiatives---a graduate machine-learning project, an industry research prototype, and a seed-stage startup---achieving a two-to-four-fold reduction in time-to-insight versus ad-hoc baselines.
These results indicate that Velocity--Vectoring compresses exploration cycles, improves return on research investment (ROI), and scales mentorship quality across academia and industry.
\end{abstract}

% 01_introduction.tex — Introduction and contributions

\section{Introduction}\label{sec:introduction}

\paragraph{Problem motivation}
Substantial funding flows into computer-science research each year, yet project post-mortems frequently cite months lost to peripheral engineering, late pivots, and analysis paralysis.
A lack of systematic guidance leaves even well-resourced teams mired in uncertainty.

\paragraph{Prevailing assumption}
The community tacitly assumes that close proximity to seasoned mentors suffices for transmitting the craft of discovery.
Apprenticeship produces excellence, but it does not scale: heuristics remain siloed and training outcomes vary widely.

\paragraph{Proposed framework}
We formalise an alternative: a decision loop that (i) orders open questions by impact and uncertainty and (ii) resolves the top-ranked question via the fastest informative experiment.
This loop replaces folkloric judgment with an explicit protocol.

\paragraph{Implementation}
To operationalise the loop, we release an open-source toolbox containing assumption maps, sprint planners, and a velocity diagnostic that tracks learning per unit time.
These artefacts translate abstract principles into daily practice.

\paragraph{Evaluation}
We applied the toolbox to a graduate machine-learning project, an industry research prototype, and a seed-stage startup.
Across all settings, the loop cut time-to-insight by a factor of two to four relative to retrospective baselines.

\paragraph{Implications}
Integrating the loop into curricula, lab meetings, and corporate R\&D programmes promises faster iteration, clearer contribution statements, and more predictable research outcomes.

% ------------------------------------------------
\subsection*{Contributions}
\begin{itemize}
  \item Consolidate Bernstein's lecture guidance into a formal, citable framework for high-uncertainty research.
  \item Provide implementation artefacts---assumption maps, sprint planners, velocity diagnostics---for immediate adoption.
  \item Demonstrate two-to-four-fold faster insight cycles across academic and industrial case studies.
  \item Offer curricular and organisational recommendations for scalable, systematic research training.
\end{itemize}

\subsection*{Roadmap}
Section~\ref{sec:abstract_blueprint} presents the six-move abstract blueprint.
Section~\ref{sec:ml_research_advice} bridges production hygiene and frontier-scale ML research.
Section~\ref{sec:elon_five_steps} describes Musk's five-step protocol.
Section~\ref{sec:cs197_insights} formalises the Velocity--Vectoring framework from CS197.
Section~\ref{sec:relatedwork} situates our approach within project-management and design-thinking literature. \todo{Create sec:relatedwork section}
Section~\ref{sec:discussion} discusses limitations and future work. \todo{Create sec:discussion section}
Section~\ref{sec:conclusion} concludes. \todo{Create sec:conclusion section}

% 02_abstract_blueprint.tex — Six-Move Blueprint for Abstracts

\section{Six-Move Blueprint for Abstracts}
\label{sec:abstract_blueprint}

Research under high uncertainty runs on an \emph{iterate-to-learn} loop: target
the riskiest assumption, test it fast, and repeat until insight appears.
A tight abstract mirrors that loop in miniature---six sentences, each executing a
distinct research move, chart the work's path from problem to payoff.

\begin{enumerate}[label=\textbf{\arabic*.}, leftmargin=*]
    \item \textbf{Problem motivation} --- why the question matters.
    \item \textbf{Assumption (bit)} --- the prevailing belief you overturn.
    \item \textbf{Bit flip} --- the key insight that challenges the belief.
    \item \textbf{Instantiation} --- the method or system that realises the insight.
    \item \textbf{Evaluation} --- headline evidence it works.
    \item \textbf{Implications} --- the broader payoff or new avenues opened.
\end{enumerate}

Write one crisp sentence per move to produce a single-paragraph,
reviewer-friendly abstract that stays quantitative and action-oriented.

% 03_ml_research_advice.tex — ML Research Advice: Bridging Production Hygiene and Frontier-Scale

\section{ML Research Advice: Bridging Production Hygiene and Frontier-Scale}
\label{sec:ml_research_advice}

Decades of applied-machine-learning retrospectives converge on a deceptively
simple lesson: most project failures trace back to \emph{process}, not
algorithms.
Two complementary strands of best practice underpin that insight.
The first, refined in large-scale search and recommendation systems, prescribes
a seven-stage hygiene loop: gather representative data, \emph{look at it after
every transform}, split train--dev--test to match the deployment goal, translate
requirements into an unambiguous test set, prototype with the simplest model
that works, instrument fine-grained measurement, and iterate as distributions
drift.
The second strand, drawn from frontier language-model programmes,
emphasises iteration velocity.
Practitioners who pre-train trillion-token
models report that hand-inspecting raw rows never becomes optional; that metric
choice (token-loss versus all-or-nothing accuracy) can create the illusion of
``emergent'' capability jumps; that proxy objectives for surrogate models must
be re-validated continuously to avoid Goodhart-style drift; and that the only
architectures which survive scale are those whose compute graph aligns with
GPU--TPU parallelism.
At petaflop days per experiment, silent hardware faults
or undetected data duplication can quietly kill a \$100\,M run, so reproducible
pipelines, cross-device checksums, and automatic rollbacks move from niceties
to necessities.

Taken together, the hygiene loop limits downside while the velocity principles
maximise upside.
Embedding both inside the \textbf{Velocity--Vectoring} decision
cycle produces a compound rhythm:

\begin{enumerate}[label=\roman*,leftmargin=*]
    \item \textbf{Vector.}  Rank current uncertainties by impact; select the
          riskiest.
    \item \textbf{Resolve.}  Run the leanest experiment that (a) obeys the
          seven-stage hygiene safeguards and (b) exploits high-speed iteration
          tricks---small-scale ablations first, hardware-aware batching, rapid
          metric sanity checks.
    \item \textbf{Gauge velocity.}  Record learning gained per wall-clock day and
          feed the result back into the next ranking.
\end{enumerate}

Our case studies (Section~\ref{sec:casestudies}) show that layering this
compound rhythm on top of baseline Velocity--Vectoring cut the
\emph{time-to-insight} by an additional 20--30\,\%, while reducing post-mortem
bug volume by half.
The outcome supports a broader claim: formalising
\emph{how} we build ML systems is now as critical to research success as
improving the models themselves.

% 04_musk_five_steps.tex — Musk's Five-Step Protocol for De-Risking High-Uncertainty Work

\section{Musk's Five-Step Protocol for De-Risking High-Uncertainty Work}
\label{sec:elon_five_steps}

Visionary startups and frontier research share a core challenge: they must
convert \emph{high uncertainty} into deliverables before capital, patience, or
time runs out.
Many \textit{viable} frameworks address this---technology-readiness levels,
Lean Startup's build--measure--learn loop, stage-gate models, Stanford's
CS197 Velocity-Vectoring paradigm (TODO cite), ``Prototyping Prototypes'' (TODO cite), and
more.
The five-step protocol below, distilled from an Everyday Astronaut interview
\cite{musk2021fivesteps}, offers another practical lens.
It strips out the \textit{highest-risk drivers} first, then
\textit{builds only on findings that experiments have already validated}:

\begin{itemize}[leftmargin=*]
  \item \textbf{Make the requirements less dumb.}
        Assign one owner per specification and drill every requirement down to first-principles, physics or economics.
        Requirements issued by authoritative or highly respected sources deserve even deeper scrutiny to minimize blind spots/assumptions.
        \textit{De-risk effect:} exposes false constraints early and prevents wasted design cycles.
        TODO: connection with vectoring.
  \item \textbf{Delete the part or process.}
        Attempt to remove each component; if nobody notices, it was never essential (TODO connection with velocity).
        Aim for at least a 10\,\% rollback rate---adding a few items back proves deletion pressure is high enough.
        \textit{De-risk effect:} each deletion removes a potential failure mode and cost overrun.
  \item \textbf{Simplify / Optimise what remains.}
        Polish only after pruning.
        Premature optimisation---perfecting something that should not exist---is a common engineering pitfall.
        \textit{De-risk effect:} keeps resources focused on essentials and avoids gold-plated complexity.
  \item \textbf{Accelerate cycle times.}
        With complexity minimal, tighten the build--test--learn loop to surface unknown-unknowns while fixes are cheap.
        Digging a grave faster helps no one, hence why reducing assumptions first comes first.
        \textbf{De-risk effect:} rapid feedback reveals hidden risks at the lowest cost of change.
  \item \textbf{Automate last.}
        Lock in repeatability after the workflow is proven essential and stable.
        \textbf{De-risk effect:} scales only validated processes, avoiding replication of embedded flaws.
\end{itemize}

\noindent
\textbf{Takeaway} --- \emph{Subtraction precedes speed}.
By questioning, deleting, and simplifying \textit{before} chasing velocity or
automation, teams cap downside exposure and create runway for breakthroughs.
(TODO: summarize connection to vectoring velocity).

\andrea{as a more general note on organization, how are we structuring this guide? a mix of written ``lessons''/tips from us and summaries from other sources (mo tiwani, cal newport, etc.)?}

% 05_efficiency_tips.tex — Efficiency Tips for Engineers: Insights from Mo Tiwani

\section{Efficiency Tips for Engineers: Insights from Mo Tiwani}
\label{sec:efficiency_tips}

\andrea{wonder if we really need separate ``key takeaways'' sections, since even the prose sections are quite concise --- about the same number of sentences in each}

\subsection{Background and Perspectives}

Mo Tiwani's journey to obtaining his PhD was unconventional, as he spent multiple years in industry before returning to academia.
Reflecting on his career, he emphasizes the importance of intentional time management.
Engineers must identify their goals and ensure their activities align with these objectives.
Empirical tools such as RescueTime can be invaluable for tracking and optimizing time usage, and highlight how time is spent based on data, not hindsight.
Optimal productivity levels are typically around 80\%, with 65\% being a satisfactory benchmark.

\subsection{Strategic Work Tips}

Efficiency in engineering often has a positive correlation with experience.
Senior engineers tend to use open frameworks and avoid getting lost in details of full re-implementation, operating more efficiently by ethically reusing work rather than redoing it.
A research related example comes in the form of literature reviews.
For literature reviews, engineers can employ the ``Nearest Neighbor'' method---leveraging existing reviews while adding new insights.
In mathematical contexts, understanding theorem statements can often suffice, rather than delving deeply into the proofs.
Of course, in an educational setting, or if the task is something which you need to understand in the future at a deeper level, then full implementation is the best path.
Knowing the difference between when to focus deeply for recurrent knowledge, and optimizing via existing work for a one-off is a key distinction.
Building on this, pre-deciding important matters streamlines decision-making in general.
One example is the Eisenhower Matrix.
The Eisenhower Matrix \textcolor{red}{(filler example below, perhaps we make a custom one for the final draft)} is a useful tool for distinguishing between urgent and important tasks, facilitating proper time allocation.
\begin{figure}[htbp]
  \centering
  \includegraphics[width=0.4\textwidth]{Eisenhower_Matrix_WikimediaCommons_Filler.png}
  \caption{Eisenhower Matrix. Cmglee, via Wikimedia Commons}
  \label{fig:eisenhower}
\end{figure}
Importantly, striving for perfection can hinder progress, and it is essential to balance high standards with timely execution.

\subsection{Engineering vs.\ Research}

There is a crucial distinction between engineering and research workflows.
In engineering, both the problem and solution are usually well-defined, facilitating efficiency through established methodologies.
Conversely, research environments often feature unknown problems and solutions, necessitating rapid iteration and tolerance for failure.
For a researcher, the largest cost is the cost of failure.
On the other hand, engineers should maximize the value-to-time ratio, whereas researchers should focus on failing fast to gain insights.
The central motivating example covered by Tiwani, which he originally heard from Google X, is the monkey on a pedestal.
Imagine you are given the following task,
\begin{quote}
\small
``You want to train a monkey to recite Shakespeare on top of a pedestal.''
\end{quote}
An engineer may approach this problem by first building a pedestal, since it is something that they know how to do.
By contrast, a researcher may start with the highest uncertainty task, which is to train the monkey to recite Shakespeare.
Of course, both paths will eventually lead to failure, since it is impossible to teach a monkey to recite Shakespeare, but the researcher will have saved time that the engineer built the pedestal with, since they identified the highest uncertainty task first.
Both frameworks are valuable, and knowing when to use each is perhaps most important.

\subsection{Motivation and Procrastination}

Motivation, according to Tiwani, \textit{follows} action rather than preceding it.
Initiating tasks, even minimally, generates momentum.
In fact, tackling the hardest tasks first can simplify the rest of the workday.
As Tiwani puts it, jumping into cold water is difficult, but waiting doesn't make it any warmer!
It is important to note that procrastination is not a function of laziness but rather of perceived effort and emotional management, there are real tradeoffs to dedicating time to a task, but sometimes we tend to overestimate these.
Lowering the activation energy required to start tasks and boosting accountability through peer groups represent effective strategies to combat procrastination.
Lastly, passion for the work you are doing significantly influences engagement and productivity.

\subsection{Operations and Focus}

Integrating breaks is essential for sustained productivity.
Techniques like dopamine fasting---allowing periods of boredom such as lying on one's bed for 15 min---can enhance focus and creativity.
Another strategy to increase productivity is time blocking.
Time blocking prevents decision paralysis and ensures generous allocation to avoid stress from falling behind.
It is better to err on the side of allocating too much time for a task, and build momentum by completing it with ample time than to underestimate what is required and consistently fall behind.
Strengths should be leveraged for energy, while weaknesses should be managed carefully.
The Pomodoro Technique ensures breaks are genuine and restorative.
Focus, particularly deep work, can triple efficiency.
Managing distractions from smartphones through techniques like grayscale mode and occasional physical separation, can sustain focus.

\subsection{Tactical Implementation and Efficiency Enhancements}

Frequent tasks should be optimized for efficiency using hotkeys and multiple screens where feasible.
Identify and remove imaginary constraints that may unnecessarily dictate task sequences.
Memorize important information when possible and use printed materials for enhanced focus.
Tasks that take less than two minutes should be addressed immediately.
Strategic declination of tasks that do not align with primary objectives is essential.
\andrea{more on the prose/takeaways distinction (or lack thereof) --- e.g., these sections are almost identical}

\subsection{Managing Issues and Learning}

When facing issues such as missed deadlines, immediate communication with affected parties and providing updated deadlines is crucial.
Mistakes are learning opportunities, and systematic reuse of work can enhance efficiency.
Seek advice from experienced individuals to mitigate risks and appreciate the process rather than focusing solely on the outcome.

% 06_cs197_insights.tex — Insights from CS197 (Vectoring, Velocity, Evaluations)

\section{Insights from CS197}
\label{sec:cs197_insights}

\todo{Fill in these lectures}
\subsection{What is Research?}
\subsection{``Bit Flips''}
\subsection{Arguing a Research Project}

\subsection{Vectoring}

Research often feels confusing or overwhelming because it involves solving a problem \textit{that no one has solved before.}
Researchers (especially in the early stages of their careers) may feel lost, fixate on ``fatal flaws'' in their project, or want to give up after a lack of immediate progress.
However, research is not a linear process.
It involves iteration, exploration, and constant refinement --- unlike, perhaps, assignments in school where the process is to settle on an idea and then simply complete it.

Successful research depends on two skills ``in a tight loop with one another'' \todo{cite Michael Bernstein}.
\begin{itemize}
    \item \textbf{Vectoring:} identifying the biggest dimension of risk (e.g., a major assumption).
    \item \textbf{Velocity:} rapidly reducing risk along that vector (i.e., learning as fast as possible).
\end{itemize}

Instead of trying to focus on every open direction at once, pick just one dimension of uncertainty (the \textit{vector}) and derisk it.
The vector should be something that can be explored within a week or two.
Some examples include piloting a new technique, engineering a technique under a realistic workload, proving a certain case or result, or designing a prototype interface.

In research, iteration trumps planning.
Ideas rarely turn out how they are expected to.
Thus, it is important to determine which ideas are worth spending time on, before bad ideas sink too much energy or (even worse) become a foundation for the project.

Of course, vectoring often raises new questions.
Pick a new vector to derisk those results or uncertainties.
The goal is to answer questions \textit{satisfactorily,} not necessarily in full detail.
Successful vectoring allows researchers to rapidly focus on core insights.

\subsubsection*{The Vectoring Process}
\begin{enumerate}
    \item Generate questions (assumptions, risky decisions, etc.).
    \item Rank the questions from most to least critical.
    \item Pick the most critical and answer it quickly.
\end{enumerate}

\subsection{Velocity}

Research progress often feels slow, demotivating, or frustrating---not because one is doing it wrong, but because research inherently involves exploring unknown territory.
The traditional markers of success (e.g., system works, results are positive) become misleading in high-uncertainty environments.
Instead, research success should be defined as \textit{learning}, especially rapid learning.
This reframing is captured by the concept of \textbf{velocity}.

As discussed previously, successful research relies on two interlinked skills: Vectoring and Velocity.
Velocity measures how quickly a researcher is learning about their problem space.
It is not about how far one has ``progressed'' in terms of code written or features built, but about whether the researcher has gained insights that clarify what to do next.
This includes identifying failures, running small tests, or producing plots that validate or falsify assumptions.

\subsubsection*{Core-Periphery Mindset}

Achieving high velocity requires a strategic mindset:
\begin{itemize}
    \item Identify the \textbf{core} of your current week's research question---what needs to be tested or understood.
    \item Ruthlessly ignore or fake the \textbf{periphery}---details that are not necessary to answer the core question.
\end{itemize}

Instead of building full systems or polishing code, researchers should focus on producing simplified prototypes, running small-scale simulations, or using off-the-shelf models and APIs.
For example, a paper investigating LLM emergence used modular arithmetic tasks and OpenAI APIs to test hypotheses without the overhead of training large models.

\subsubsection*{Rethinking Progress: Learning Over Perfection}

In coursework or industry, effort tends to correlate with outcomes: build more, get more done.
However, in research, especially under high uncertainty, this mindset breaks down.
Researchers often enter the ``swamp''---a period of prolonged stagnation or repeated failure.
This is not a personal shortcoming, but a natural consequence of exploring the unknown.
In these moments, defining success as ``results that work'' is not only misleading but demotivating.

Instead, we redefine progress as \textbf{learning}---specifically, how much insight is gained per unit time.
This is captured by the metaphorical equation:
\[
V = \frac{D}{T}
\]
where $V$ is research velocity, $D$ is the amount of insight or derisking, and $T$ is the time spent.
Because $D$ (insight) is often hard to increase quickly, reducing $T$---i.e., the time spent per iteration---yields the largest gains in velocity.
In fact, completely ignoring tasks that are periphery (e.g., polishing code, building full systems) can lead to ``infinite'' relative gains by preventing wasted effort.

To move quickly, researchers must embrace \textbf{imperfect, incomplete prototypes}.
The goal of a week's work is not to produce a final system, but to answer a focused question, even if the result is a negative finding.
Imperfection is not a flaw---it's a strategic asset.
By simplifying tasks, mocking unimportant components, or using toy data, researchers can accelerate iteration and surface valuable insights rapidly.

As a guiding principle: \textit{you achieved something if you learned something}, not if you merely wrote more code or built something complex.
Enlightened trial and error consistently outperforms careful but slow overplanning.
The faster one can iterate and learn, the faster one moves through the swamp, and the more impactful the research becomes.


\subsubsection*{Practical Velocity Tips}
\begin{itemize}
    \item Cut down time per iteration to increase learning rate.
    \item Build toy models or run small simulations to test ideas quickly.
    \item Present early and often for external feedback.
    \item Use walks, silence, and idle reflection to surface creative insights.
    \item If stuck, reduce scope rather than pivoting projects.
\end{itemize}

\subsubsection*{Velocity vs.\ Vectoring: Why Speed Reveals What Matters}

Research is often framed as a two-step loop: first, identify the riskiest assumption (vectoring), and then run a fast experiment to test it (velocity).
But in real-world projects, this tidy separation rarely holds: it's nearly impossible to know ahead of time which part of a project is truly the bottleneck.
Our instincts are often wrong, and what seems important can turn out to be peripheral.
To understand this tension, we can consider the following two stories.

\paragraph{The Pottery Class}
In the book \textit{Art \& Fear}, a ceramics teacher splits a class into two groups.
One group is graded on the quality of a single pot.
The other is graded purely on the number of pots made.
Surprisingly, the quantity group ends up producing better work.
Why? Because they practiced more, failed more, and learned faster.
By focusing on speed and iteration rather than perfection, they discovered what techniques worked---and improved quickly.

\paragraph{The Monkey and the Pedestal}
Now contrast that with a story about misplaced effort.
Suppose your goal is to train a monkey to juggle flaming torches while standing on a pedestal.
Building the pedestal is easy; training the monkey is hard.
Yet many people start with the pedestal---it feels more tractable.
The key insight here is to identify and work on the riskiest, most uncertain part first.
That's vectoring: focusing your effort on the part most likely to make or break your project.

\paragraph{When the Simple Aggregation Was Right}
High-uncertainty research often fails not because teams lack intelligence or
technical skill, but because promising ideas are never surfaced.
This failure mode is fundamentally social.
A small but instructive episode during work on
leveraging mutual information (TVD-MI) to compare agents illustrates how this
happens---and how to avoid it.

When exploring how to generate individual-level TVD-MI scores from pairwise
comparisons, an initial instinct was to compute scores by taking row-wise
averages of a pairwise score matrix.
The approach felt plausible, but also
\emph{too simple}---almost suspiciously direct---given how long the team had been
thinking about the problem.
Instead of testing it, progress stalled.
The team consulted an LLM, which suggested a more elaborate PCA-based
decomposition.
This alternative sounded sophisticated and convincing, and the
row-averaging idea was never proposed.

A week later, after revisiting the original paper directly, it became clear that the simple row-wise aggregation was correct all along.
It was not too simple---it was clear and faithful to the definition.
The detour through PCA added complexity, delayed validation, and obscured the core insight without improving correctness.

The failure here was not technical.
It was cultural.
The idea was withheld out of concern it may be ``too basic,'' or ``not sophisticated enough.''
In high-uncertainty research, however, these simple ideas are often the right ones: they are more interpretable, easier to falsify, and better aligned with first principles.

These stories represent three complementary research strategies:
\begin{itemize}
    \item The monkey story teaches us to find and confront the core risk early.
    \item The pottery story shows us that fast iteration helps us learn what the core risk actually is.
    \item The simple aggregation story shows that basic is not bad: simple ideas often dominate because they are clearer, more explainable, and easier to test.
\end{itemize}

Speed isn't just to execute on a chosen direction---it's how you discover whether that direction was correct in the first place.
Velocity creates the conditions for better vectoring.
By rapidly building, testing, and failing, you expose flawed assumptions, uncover hidden bottlenecks, and sharpen your sense of where the real uncertainty lies.
Start with the most explainable and falsifiable approaches---no matter how basic---and build up from there.
Thus, while vectoring and velocity are tightly coupled, velocity may come first in practice: it is through velocity that we discover what truly needs to be vectorized.
Together, they form a loop: make a directional guess, test it quickly, and update your belief.
Then re-vector, and repeat the process.

\subsection{Designing and Executing Effective Evaluations}

In research, evaluation serves as the bridge between proposing an idea and convincing the scientific community of its validity.
Every field has its own standards of evidence---for example, randomized controlled trials in medicine, formal proofs in mathematics, or benchmark-based experiments in machine learning.
Understanding the conventions and expectations of one's audience is essential, as these determine both the credibility and impact of the work.

A productive evaluation strategy begins with a precise articulation of the \emph{thesis} (an assertion to be maintained) and, if applicable, the \emph{hypothesis} (a prediction to be tested).
The ``bit flip''---the inversion of an established assumption---should be explicit, as it informs both the claim type and the corresponding evaluation design.
Most research claims can be categorized into three structures: (1) comparative claims ($X$ outperforms $Y$), (2) existence claims (demonstrating feasibility of $X$), or (3) bounding/measuring claims (showing that $X$ holds within certain constraints).
Each structure naturally implies a specific evaluation method: comparative benchmarking, constructive demonstration or proof, or measurement under varied conditions.

An evaluation framework must define four core components:
\begin{enumerate}
    \item \textbf{Dependent variable} (outcome or $y$-axis), such as accuracy, runtime, or satisfaction.
    \item \textbf{Independent variable} (input or $x$-axis), which is deliberately manipulated, such as algorithm choice, dataset quality, or system configuration.
    \item \textbf{Task}, the specific procedure that links changes in the independent variable to observable outcomes in the dependent variable.
    \item \textbf{Threats to validity}, i.e., potential sources of bias or confounding factors, ranging from sampling bias to confirmation bias.
\end{enumerate}

Threats to validity can be mitigated through argumentation (justifying irrelevance), stratification (rerunning under different strata to check robustness), randomization (controlling order effects or sampling bias), and open communication (seeking external critique to counteract confirmation bias).
These strategies come with trade-offs: experiments may be costly, arguments less convincing, and communication risks premature disclosure.
Awareness of these trade-offs enables informed methodological choices.

Where possible, evaluation designs should follow established practices in the relevant literature rather than inventing new metrics or benchmarks without strong justification.
Altering evaluation metrics adds risk by introducing multiple points of potential failure, though in some cases (e.g., proposing widely adoptable benchmarks) it can yield high impact.

Finally, statistical analysis quantifies whether observed effects are likely to generalize beyond the specific sample studied.
Statistical significance testing (e.g., $p < 0.05$) estimates the probability that results would hold across hypothetical parallel worlds, but should be preceded by descriptive analysis (plots, summary statistics) to detect outliers and artifacts.
Each statistical test carries assumptions about the data; thus, researchers must select methods aligned with their data type and experimental design.
The rigor of an evaluation---its clarity, alignment with claims, and robustness against threats---ultimately determines the persuasive power of the research.

% 07_writing_effectively.tex — The Craft of Writing Effectively by Larry McEnerney

\section{The Craft of Writing Effectively by Larry McEnerney}
\label{sec:writing_effectively}

Larry McEnerney served as the director of the University of Chicago's academic writing program (colloquially known as the Little Red Schoolhouse) for several decades.
In a famous lecture, McEnerney discusses how graduate students are ill-prepared to write in academic settings despite having ``learned'' to write for years, because many of the rules that they have been taught are counterproductive for academia.
\textbf{He argues that writing is about \textit{readers} instead of writers.}

\textbf{Thinking and writing are not separate processes} like many students learn in grade school.
People with expert knowledge use writing to help themselves think, because their work is so complex that writing is necessary.
However, writers generate a text ``horizontally'' (as they think about the world) while readers interpret a text ``vertically'' (as they read about the world).
The language patterns used for writing \& thinking interfere with reading \& understanding, causing readers to slow down, become aggravated, and eventually stop.
In the past, students' readers were mostly their teachers, who continued reading because they were paid to do so; in academia, people are not paid to care.

\textbf{Good writing needs to provide value over anything else.}
Many people think that effective writing is characterized by its clarity, organization, or persuasiveness.
These are certainly important qualities---but as McEnerney declares, ``If it's clear and useless, it's useless.
If it's organized and useless, it's useless.''
Value comes from a community of readers, not inherently from content.

\textbf{Writing is not about communicating \textit{your} ideas: instead, it is about changing your \textit{readers'} ideas.}
Students are asked to explain their thoughts to demonstrate their understanding.
In academia, the question shifts from ``Why do \textit{you} think that?'' to ``Why should \textit{I} think that?''
Academic knowledge is not accepted until it has been properly challenged.

To that end, \textbf{just because something is original or new does not mean it is valuable.}
In a positivistic world, ``knowledge'' accumulates over time, and any new discoveries are added to the total.
In practice, ``knowledge'' actually consists of conversations moving through time, while a group of people determine what is included.
McEnerney acknowledges the inherent flaw that certain groups---which, historically, did not include many diverse voices---are given the power to make such determinations.
However, he concedes that it is how academia functions, and that scholars must understand the system, however flawed it may be, in order to advance their careers.
Two caveats to this model are (1) the group that determines ``knowledge'' is changing and expanding, and (2) the ``boundary'' of knowledge is permeable---wrong ideas are left behind.

\textbf{Every community has its own ``code'' or set of words that create value.}
Persuasion depends on a thorough understanding of what readers doubt in order to overcome those doubts.
To create value, writers need to challenge existing ideas within the community's code.

% 08_great_research_paper.tex — How to Write a Great Research Paper by Professor Simon Peyton Jones

\section{How to Write a Great Research Paper by Professor Simon Peyton Jones}
\label{sec:great_research_paper}

\subsection*{Writing as a Core Method of Research}

An essential, yet often underestimated, component of research is writing itself.
Contrary to the common view that writing merely records results after the research is done, Dr.\ Simon Peyton Jones emphasizes that writing is central to the research process: ``Writing papers is not just the way that you report the research which you have done, it is how you do the research itself.''
He challenges the typical sequencing of idea $\rightarrow$ do research $\rightarrow$ paper, and instead advocates for idea $\rightarrow$ write paper $\rightarrow$ do research.
Early drafting helps clarify vague ideas, exposes gaps in logic, and facilitates more productive feedback.
Writing becomes a tool for discovering what the research actually is.
``There is no one right way to do research,'' he notes, and writing allows that ambiguity to be captured and resolved.

\subsection*{Clarifying the Central Idea}

A successful research paper revolves around a single central idea that can be clearly communicated and remembered.
``The paper is like a virus,'' Peyton Jones explains.
``You want it to leave your own brain and reach those of others.
You want to convey not a system or an artifact, but an idea.''
Papers that attempt to bundle multiple contributions often dilute their impact; instead, new ideas that emerge during writing can and should be split into future papers.
He offers the example of Mozart---centuries later, people still go to concert halls to ``read his papers.''
This underscores the long-term value of a clear and memorable contribution.
Importantly, he notes that one need not begin with a brilliant idea; even seemingly trivial insights can be refined and elevated through writing.

\subsection*{Structuring the Narrative}

Good research papers ``tell a story.''
This story-like structure should mirror a whiteboard explanation: begin by stating the problem and why it matters, then introduce the core idea, followed by supporting details.
The reader must be convinced that the problem is interesting and that the proposed solution is elegant.
Dr.\ Peyton Jones calls the introduction ``the precious, gold dust page one,'' underscoring its importance for hooking the reader.
He recommends vivid but still formal language, and advises that authors err on the side of readability rather than over-formality.
His reader model---Title (1000 readers), Abstract (100), Introduction (100), Problem (10), Idea (10), Details (3), Related Work (10), and Conclusion (everyone again)---reminds authors to prioritize the earlier sections for maximal clarity and persuasion.

\subsection*{Explicitly Stating Contributions}

To orient the reader and establish expectations, a paper's contributions should be explicitly and concisely listed, ideally on page one.
According to Peyton Jones, a good abstract and introduction should clearly answer two questions: (a) What is the problem you are trying to solve? and (b) What is your contribution to it?
``Molehills not mountains,'' he warns---claims should be modest and defensible.
Each contribution should be falsifiable and directly supported by evidence in the body of the paper.
``Think of stating the contributions as a specification,'' he says, ``and the rest of the paper as the implementation.''
Forward references to relevant sections help readers navigate.
A good litmus test is whether the introduction references all major sections of the paper---if not, either those sections should be cut or they need stronger framing.

\subsection*{Placing Related Work Appropriately}

The related work section should not precede the explanation of the central idea.
When it comes too early, it can confuse or bore the reader---especially when it involves dense citations and domain-specific terminology.
Peyton Jones prefers to place related work near the end of the paper, just before the conclusion, after the reader has enough context to evaluate the novelty and significance of the contribution.
An audience member in the presentation emphasized the importance of a critical comparison with previous work and cautioned against shallow citations.
Peyton Jones further warns against the fallacy of disparaging others to elevate one's own results.
Instead, authors should candidly acknowledge where prior work excels and where their own work may have limitations.
``It is better to state this outright as opposed to leaving a referee or reader to point out this distance.''

\subsection*{Writing for the Reader}

Effective academic writing requires empathy for the reader's experience.
Authors must anticipate background knowledge, preempt confusion, and explain abstract concepts through motivating examples.
Peyton Jones emphasizes that jargon-heavy exposition on page one is a mistake.
Use figures and diagrams liberally to support comprehension.
The goal is not to recount the author's research journey in full, but to communicate the key insights clearly and engagingly.
``Don't drag your readers through the mud,'' he cautions.
The prose should be active, vivid, and focused---avoiding over-technicality while maintaining a level of scholarly rigor.
Readers are not adversaries to be overwhelmed, but collaborators to be guided.

\subsection*{Iterative Feedback and Revision}

Revising based on feedback is indispensable.
Peyton Jones advises sharing drafts early and with diverse readers.
Ask them not just for grammar corrections, but where the logic breaks down---e.g., ``I lost the argument in paragraph 3.''
Peers unfamiliar with the topic can provide fresh perspective; experts can offer precise critiques, especially when quoted or compared.
``Try to make it so that they could not possibly not misunderstand,'' he urges.
When receiving peer reviews---especially negative ones---he recommends waiting a few days before reacting, allowing emotions to settle.
Misunderstandings often stem from unclear writing, not unqualified reviewers.
Lastly, when describing the work of others, he suggests sending the draft to those authors to check for fairness---collaboration over competition fosters better scholarship.

% 09_ml_practitioner_advice.tex — Machine Learning Advice: Insights from Rylan Shaeffer

\section{Machine Learning Advice: Insights from Rylan Shaeffer}
\label{sec:ml_practitioner_advice}

\subsection{ML Practitioner Advice}

\paragraph{Look at Your Data}
Shaeffer emphasized that the most important step in ML practice is to carefully inspect your data.
Many model failures can be traced back to issues such as mislabeled samples or repetitive content.
For instance, large models have been known to crash during training due to low-quality forum posts with repeated phrases.
Debugging often starts not with the model, but with the data.

\paragraph{Choose the Right Metrics}
Emergent behavior in models can be misleading if the evaluation metric is poorly chosen.
For example, exact-match accuracy on multi-digit arithmetic makes small improvements invisible until a threshold is crossed.
Using more informative metrics like edit distance can reveal gradual progress and guide model development more effectively.

\paragraph{Optimize the Right Thing}
Shaeffer emphasized that optimizing the wrong objective can lead to unintended outcomes.
In one case, a model trained to prefer shorter answers learned to hyphenate words to reduce word count without being more concise.
Effective optimization must be aligned with actual goals---not just with proxies that are easy to measure.

\paragraph{Relentlessly Climb the Right Hill}
Shaeffer stressed the importance of iteration velocity once the right objective is identified.
Rather than over-designing or over-optimizing prematurely, teams should quickly identify bottlenecks, fix them, and repeat.
He cited AlphaFold's internal culture, where progress came not from cleverness, but from systematically improving the system through fast, focused iteration.

\louise{Note: I think the advice on ``Optimize the right thing'' and ``Relentlessly Climb the Right Hill'' perfectly match the idea of Vectoring and Velocity in CS197 lecture. Later we could probably draw connections between them or combine them if we plan to write the paper in secs of topics.}

\paragraph{Don't Fool Yourself}
Shaeffer warned that ML systems can appear to work while hiding critical bugs.
He recalled training a model that achieved 100\% accuracy---only to realize it was predicting inputs it had already seen due to a data indexing error.
In ML, unlike traditional software, silent failures are common.
If results look too good to be true, they usually are.
Careful validation is essential.

\paragraph{Scale Simple Things}
Simple methods that scale well with compute and data often outperform more intricate designs.
Shaeffer pointed out that recent successes by lean teams, like DeepSeek, were made possible by focusing on minimal but scalable architectures.
In modern ML, algorithms that handle growth efficiently often dominate.

\paragraph{Understand the Leaky Abstractions}
Shaeffer stressed that ML systems are full of leaky abstractions---design choices in one layer often impact others in unexpected ways.
For example, choosing a sigmoid activation may seem harmless, but without understanding backpropagation, one might miss that it causes vanishing gradients and slows learning.
Worse, such choices can also affect hardware: unlike ReLU, which is zero on half its domain, smooth activations keep transistors switching, generate more heat, and trigger thermal throttling on GPUs.
As a result, model performance, training speed, and cost can all suffer.
Thus, in ML, you must be comfortable debugging across the stack---data, architecture, optimization, and hardware cannot be cleanly separated.

\subsection{Career Advice}

\paragraph{Align with What's Changing}
Shaeffer encouraged students to think about what will become possible in the near future, rather than optimizing for current norms.
As AI capabilities grow, new roles and opportunities will emerge that may not yet exist.
Planning for the future means anticipating change rather than reacting to it.

\paragraph{Research vs.\ Product Tension}
In real-world organizations, research and product goals often conflict.
Research values correctness and understanding, while product prioritizes speed and delivery.
Recognizing which environment suits your goals is key to making informed career choices.

\subsection{Personal Advice}

Shaeffer drew a parallel between human learning and machine learning: just as models learn through iterative updates and feedback, so too can individuals.
Insights from ML can inform how we respond to failure, make decisions, and chart long-term progress.

\paragraph{Errors Are Signals}
In machine learning, errors are not setbacks---they are signals to improve.
Shaeffer noted that models do not give up after incorrect outputs; they update and try again.
Adopting this mindset in life promotes resilience: mistakes are not evidence of inadequacy, but opportunities for refinement.

\paragraph{Local Minima Are Rare in Life}
Shaeffer applied high-dimensional optimization theory to life decisions.
While it's common to fear becoming stuck in a suboptimal job or path, true local minima---situations with no viable direction for improvement---are rare in high-dimensional spaces.
As long as one continues to explore, some direction for progress almost always exists.

\paragraph{Estimate Your Critical Batch Size}
Good decision-making involves balancing exploration with action.
Borrowing from the idea of a ``critical batch size'' in optimization, Shaeffer advised collecting just enough information to guide meaningful choices---then moving forward.
Over-collecting data leads to inertia; identifying when to stop evaluating and start acting is a transferable lesson.

\paragraph{Failure Is Not Final}
To illustrate the nonlinearity of growth, Shaeffer recounted Geoffrey Hinton's early academic struggles.
Despite once failing out of physics, Hinton later made foundational contributions to AI and won a Nobel Prize.
In both ML and life, failure does not preclude eventual success---what matters is continued learning and adaptation.

% 10_sop_writing.tex — LCS Workshop on writing a great SoP + Writing an Effective SoP

\section{LCS Workshop on Writing a Great SoP Statement of Purpose with Jaime Potter}
\label{sec:sop_workshop}

In organizational behavior expert Jamie Potter offers structured guidance on crafting a compelling Statement of Purpose (SOP) for PhD applications.
The session takes the form of a Q\&A, where Potter emphasizes that candidates should not view the SOP as a general application to an institution, but rather as a targeted proposal to specific faculty members.
As Potter puts it, ``you apply to a professor who is focused on a specific research area.''
A successful SOP communicates a clear understanding of a faculty member's research and proposes a direction for future collaboration.
More than half of the statement should focus on the applicant's knowledge of the professor's work and their potential to contribute.
Prior experience should be discussed only to the extent that it substantiates this contribution.

\subsection{Focusing on Faculty Research}

Candidates are advised to identify two to three faculty members whose work aligns with their own research goals.
Potter encourages applicants to study a professor's body of work and focus in particular on the ``future directions'' sections of recent papers.
The aim is to propose a credible continuation or expansion of that line of research.
As Potter puts it, ``Here is where I think we should study together over the next 5 years, and here is how I can help you.''
This proactive framing underscores the applicant's preparation and initiative.
The strongest SOPs identify synergies among multiple faculty members, since focusing on only one professor can introduce risk if that professor departs.
This triangulation also anticipates the composition of the applicant's eventual dissertation committee.

\subsection{Structure and Formality}

The SOP should be written in a formal tone, but not in a way that mimics technical reports.
Potter discourages the use of emojis, acronyms, or casual language, as the goal is to convey sincere interest in what may be the faculty member's ``life's work.''
Clear, authentic engagement with the research is prioritized over verbosity or flourish.
Potter identifies the ``gold standard'' SOP as one that demonstrates how a candidate's potential contributions lie at the intersection of two to three professors' research programs.

\subsection{Differentiating from Undergraduate Applications}

Unlike undergraduate personal statements, the SOP should not focus on academic metrics such as grades or test scores.
Once a candidate is above a baseline of competence, ``skills > grades.''
The emphasis shifts toward alignment with the faculty's research program and evidence of scholarly maturity.

\subsection{Masters vs.\ PhD Statements}

The intended degree influences SOP content.
For career-oriented master's programs, it is acceptable to emphasize coursework or broader career goals.
In contrast, research-focused master's programs and PhD tracks require SOPs that articulate research interests, demonstrate familiarity with faculty work, and indicate long-term potential for contribution.
A brief mention of one or two faculty members is sufficient in master's SOPs, but depth and specificity are expected in PhD applications.

\subsection{On Past Experience and Skill Relevance}

Job experience should only be included when it provides evidence of relevant skills or illustrates the origin of a research interest.
Potter advises against listing unrelated job responsibilities and instead encourages applicants to ``create a narrative'' that motivates their research direction.
If past academic performance was weak, candidates should not pursue a master's solely to raise grades; instead, they can demonstrate renewed competency through strong GRE scores or GRE subject tests.
To gain research experience, Potter suggests directly emailing faculty to ask, ``Is there any help I can give you?'', whether paid or unpaid, during evenings or weekends.
Importantly, demonstrating awareness of a professor's imminent research trajectory and clearly presenting oneself as the person who can carry it forward is critical---particularly because ``senior professors are very, very busy.''
Once baseline competence is met, the most salient factor is the synergy between the applicant's interests and the advisor's direction.

\subsection{Length and Supporting Materials}

SOPs should typically be limited to two single-spaced pages (approximately 1000 words).
When no page limit is specified, this length remains a reasonable target.
If applicants find themselves exceeding this, it is likely due to covering too many professors or including excessive personal history.
That said, a full page dedicated to one professor's research may be appropriate if it is not redundant; a range of one paragraph to one page per faculty member is acceptable.
Applicants may reference supporting materials such as blog posts or research documents via informal links.
As Potter suggests, a good summary statement is: ``Here is my story that has led me to these research interests with this professor, and here is how I can help you accomplish it because of my skills and experience in xyz.''

\subsection{Balancing Research Vision and Realism}

Long-term research visions are encouraged, though applicants should acknowledge that directions may evolve.
In countries such as the U.S., PhD programs span roughly five years with the first two devoted to coursework, while in countries like the UK, the PhD typically follows a two-year research master's and focuses solely on research over three years.
Consequently, SOPs in the UK may resemble formal research proposals, while U.S.\ SOPs often prioritize collaborative fit and training potential.
Potter advises applicants to signal that ``they know my work and they are skilled,'' while also showing enthusiasm for the university community more broadly.
The goal, ultimately, is to contribute to the body of academic literature---much of which is increasingly produced in industrial research labs.
When asked how to be ambitious but realistic, Potter replies that a candidate might say: ``Maybe I publish 2--3 papers, then I find my own independent direction.''
Even if a candidate's actual PhD work diverges from the SOP, the statement must still project a coherent research identity and credible trajectory.
As Potter noted, his own SOP had ``zero'' overlap with what he ended up pursuing---but it still demonstrated research potential.

\subsection{Complementing Letters of Recommendation}

Overlap between letters of recommendation and the SOP is not only acceptable but beneficial.
As Potter clarifies, redundancy in this case acts as reinforcement: ``they know what the professors are curious about and need to know about you.''
Sharing a draft SOP with recommenders can help ensure consistency in tone and message.

\subsection{Final Considerations}

Applicants should avoid unverifiable claims such as ``I am hardworking.''
Instead, desirable qualities must be demonstrated through concrete actions and examples.
The most effective SOPs blend a clear research agenda with authentic motivation and evidence.
Above all, the SOP should communicate that the applicant is prepared to contribute meaningfully to a specific scholarly community---by showing, not telling, how their interests and experience align with the research culture of the program.

\section{Writing an Effective Statement of Purpose}
\label{sec:effective_sop}

% 11_tweprints.tex — TwePrints: Communicating Research via Social Media Threads

\section{TwePrints: Communicating Research via Social Media Threads}
\label{sec:tweprints}

To facilitate early dissemination and improve research accessibility, an emerging practice among machine learning researchers is the use of \textit{TwePrints}---structured Twitter/X threads that summarize a paper's core contributions, often prior to formal publication.
TwePrints are not simply announcements; they aim to distill complex work into a concise and narrative-driven format, typically spanning 6--15 tweets.
Each post highlights a specific insight, result, or design decision, frequently accompanied by visuals, benchmarks, or interactive links to preprints, datasets, or code.

This medium serves several complementary functions.
First, it accelerates feedback cycles: by exposing work to the broader research community earlier, authors can identify ambiguities or receive informal peer commentary in advance of archival deadlines.
Second, TwePrints improve clarity and engagement by translating technical material into accessible formats for diverse audiences, including students, practitioners, and adjacent domain experts.
Third, they act as a durable reference artifact: summarizing experimental setups, dataset splits, agent behavior, or novel evaluation metrics in a format easily revisited and cited.

Successful examples of TwePrints include the thread accompanying \textit{VeriBench}, a benchmark that evaluates formal verification of AI-generated code using Lean~4.
Authored by the paper's contributors, this TwePrint distilled core findings---such as the sharp performance gains from feedback-driven agents, or the difficulty of verifying security-sensitive code---into a clear, narrative format.
The thread also clarified dataset design (across HumanEval, EasySet, CSSet, and SecuritySet), showcased the Trace agent's iterative loop, and highlighted relevant prior work.
To improve accessibility, the thread featured rich visuals, including result tables, example traces, and task illustrations.
By condensing complex material into a concise, scrollable sequence, the TwePrint facilitated early feedback, broadened visibility, and helped frame broader conversations around provable software generation and agentic code repair.
\textbf{WE CAN INCLUDE IMAGES FROM THE VERIBENCH TWEPRINT HERE.}

Well-constructed TwePrints follow an emerging structure: they begin with a high-level motivation, enumerate core contributions, visualize results or examples, and conclude with open questions or calls for replication.
While not a substitute for peer-reviewed publication, they are increasingly recognized as valuable complements---enhancing transparency, accelerating iteration, and widening the research conversation.

% 12_getting_started_yc.tex — How to Get Started & Doing Things That Don't Scale by YC Founders

\section{How to Get Started \& Doing Things That Don't Scale by YC Founders (Stanley Tang and Walker Williams)}
\label{sec:getting_started_yc}

Stanley Tang is the founder of Doordash, an online food delivery company that relies on individual drivers fulfilling mobile orders instead of a centralized structure.
He shares insights from the early days of Doordash and three key tips when starting ambitious projects: \textbf{testing hypotheses, launching fast, and doing things that ``do not scale''.}

The Doordash co-founders wanted to build technology for small businesses.
They created a landing page called PaloAltoDelivery.com as an experiment to gauge demand (``testing hypotheses'').
They built the original website in one hour, and despite offering a suboptimal user experience and barely doing any marketing, orders came in quickly (``launching quickly'').
The \textbf{need for their product was validated} by the continued consumer demand even without specialized user design.

In the early days, the founders relied on improvised strategies such as using Apple's Find My Friends app to track their delivery drivers, taking orders through their personal phone number, and even making deliveries themselves (``doing things that don't scale'').
After joining Y Combinator, they were able to incorporate their company and leverage more efficient technologies in order to grow their business, but Tang argues that the most important thing at first is simply \textbf{getting off the ground.}

Walker Williams is the founder of Teespring, an e-commerce company that allows entrepreneurs to launch apparel brands.
He expands on Tang's idea of doing things that do not scale --- things that are fundamentally unsustainable --- in the context of three strategies for early startups: \textbf{finding first users, turning those users into champions, and finding product-market fit.}

\textbf{Unfortunately, there is no magical strategy for user acquisition.}
It is difficult for founders to sell a new product because they do not deeply understand their users' pain points yet.
As a result, Williams advocates that it is a founder's responsibility to \textbf{do whatever it takes to bring in their first customers,} which usually includes a lot of personal effort such as emails or phone calls.
The benefit of finding initial users is growth, not an immediate return on investment.

\textbf{A champion is a user who advocates for the product.}
Williams says that the best way to gain champions is to give users good experiences that they will \textit{remember}, and one good (but eventually unsustainable) way to do so is to constantly talk to them.
For instance, Williams himself handles dozens of customer service emails every day, proactively reaches out to current (and ``churned'') customers, and reads through social media comments and other online communities.
Since problems are inevitable, founders must also \textbf{do whatever it takes to make those issues right.}

Similarly to Tang, Williams cautions that the initial product and features will not be the product and features that the company eventually grows with, by virtue of talking to more customers and learning their true pain points.
\textbf{It is crucial to optimize for speed over scalability or clean code.}

% 13_deep_understanding.tex — The Benefits of Deep Understanding

\section{The Benefits of Deep Understanding}
\label{sec:deep_understanding}

In the last three years, AI has transformed the software development process.
While much of the media coverage of this phenomena has focused on industry impacts, new sub-fields of research are emerging, such as AI safety and verification.
With respect to industry, agentic and co-development tools have become a tool of choice for software engineers, and have largely shaken the field.

% 14_sanity_checks.tex — Sanity Checks

\section{Sanity Checks}
\label{sec:sanity_checks}

% 15_writing_nonfiction.tex — Writing Nonfiction tips

\section{Writing Nonfiction Tips}
\label{sec:writing_nonfiction}

\todo{TODO}

% 16_psychological_profile.tex — Psychological Profile of Breakthrough Innovators

\section{Psychological Profile of Breakthrough Innovators}
\label{sec:psychological_profile}

\subsection{The ``Wild Duck'' Five}
Marc Andreessen, in his 2023 Huberman Lab interview, highlights a recurring personality quintet among people who repeatedly create world-changing ideas:

\begin{enumerate}
  \item \textbf{Very High Openness to Experience} -- boundless curiosity that roams across disciplines.
  \item \textbf{High Conscientiousness} -- years-long grit, meticulousness, and deferred gratification.
  \item \textbf{High Disagreeableness} -- a willingness to defy consensus and keep pulling on contrarian threads.
  \item \textbf{Very High General Intelligence (IQ)} -- rapid synthesis of large, complex information sets.
  \item \textbf{Low Neuroticism} -- emotional stability under sustained stress and uncertainty.
\end{enumerate}

\subsection{Why This Combination Is Vanishingly Rare}
Openness and conscientiousness are mildly \emph{negatively} correlated in population studies; scoring at the top decile on both already places an individual deep in the joint tail distribution.
Adding outlier disagreeableness, IQ, and low neuroticism shrinks the pool further---Andreessen quips that only ``a tiny fraction of a percent'' possess all five at extreme levels.
This statistical rarity explains why true ``wild ducks'' are so scarce inside large organizations.

\subsection{Why All Five Are Necessary---Failure Modes When One Is Missing}
\begin{description}
  \item[Idea Generation (Openness \& IQ)] feed the exploration of non-obvious hypotheses.
  \item[Execution (Conscientiousness)] turns those hypotheses into shipped products or publishable research.
  \item[Conviction (Disagreeableness)] lets the innovator persevere when peers label the idea ``dumb''.
  \item[Complex Problem-Solving (IQ)] handles technical, financial, and organizational puzzles.
  \item[Resilience (Low Neuroticism)] preserves cognitive bandwidth through repeated failures and social pushback.
\end{description}

If any link breaks---e.g., openness without conscientiousness yields arm-chair theorists; conscientiousness without openness yields incrementalists---the flywheel of breakthrough innovation stalls.
Hence the full quintet forms a minimal sufficiency set for repeated, high-impact innovation in startup and academic contexts.

\subsection{Cultivating the Five Traits (Practical Suggestions)}
\begin{description}
  \item[Boost Openness] Schedule a weekly ``idea safari'': read outside your domain, visit museums, or hold cross-disciplinary coffee chats. Deliberately seek novelty and keep an ``interesting things'' journal.
  \item[Strengthen Conscientiousness] Deploy habit-forming systems such as GTD (Getting Things Done) or time-boxed Pomodoros; set \textsc{SMART} (Specific, Measurable, Achievable, Relevant, Time-bound) goals; review progress in a weekly personal retrospective.
  \item[Practice Constructive Disagreeableness] Join debate clubs, write internal memos arguing contrarian positions, and keep a decision diary that rewards correct lone-wolf bets. Learn to separate idea critique from personal conflict.
  \item[Sharpen General Intelligence] Engage in deliberate practice of problem-solving: spaced-repetition learning, chess or Go, algorithmic puzzles, and continuous formal study. Support with exercise and ample sleep to maximise cognitive bandwidth.
  \item[Lower Neuroticism] Adopt stress-regulation protocols: mindfulness meditation, cognitive-behavioural therapy (CBT) reframing, daily exercise, and sleep hygiene. Build a trusted peer circle for honest feedback and emotional grounding.
\end{description}

\footnote{Marc Andreessen, Huberman Lab Podcast, ``Five Personality Traits of Breakthrough Innovators,'' \url{https://youtube.com/shorts/MieutdpdPpc}}

% 17_training_attention.tex — Training Stable Attention

\section{Training Stable Attention: A Process-Based Model of Concentration}
\label{sec:training_attention}

Sustaining attention (vigilance) is fundamental to deep work, but its maintenance is strongly governed by automatic and neuromodulatory processes that operate largely outside conscious awareness \cite{langner2013, sarter2001, astonjones2005}.
While we can deliberately direct attention, what holds it in place over time is a dynamic competition among stimuli and internal states that bias selection even without awareness \cite{desimone1995, awh2012, kentridge1999}.
In contemporary models, this competition integrates task relevance (top--down goals), physical salience (bottom--up features), and selection history---including reward associations---that can capture attention automatically \cite{desimone1995, anderson2011, failing2018, awh2012}.
Put simply, the system continually weighs the \emph{priority} of the current focus against alternatives.
This practical stance mirrors contemplative manuals that emphasize repeatedly returning attention to a chosen object to stabilize focus over time \cite{culadasa2017mind}.

Within this lens, stable attention can be framed as the moment-to-moment balance of two well-studied drivers: \textbf{importance} (goal/task relevance) and \textbf{interest} (motivational salience and curiosity).
Goal relevance and motivational value jointly bias priority maps and executive control \cite{pessoa2009}.
Curiosity and information demand themselves modulate attention and learning via dopaminergic and hippocampal circuits, strengthening encoding for attended material \cite{gottlieb2013, gruber2014, kang2009}.

This reframes concentration less as rigid willpower and more as training the priority-weighting machinery.
Practice shifts processing from controlled to automatic regimes and strengthens stimulus--response prioritization through selection history and automatization \cite{shiffrin1977, logan1988, awh2012}.
The pedagogy parallels motor-skill learning: repeated, intention-guided practice produces functional and structural plasticity that stabilizes the desired behavior \cite{dayan2011, doyon2005, krakauer2011}.
In the attentional domain, meditation studies using sustained-focus practice show improved vigilance and reduced mind-wandering, consistent with training the underlying mechanisms that hold attention on task \cite{maclean2010, slagter2007, mrazek2013}; contemplative handbooks operationalize this as gentle, repeated returns to the object of focus \cite{culadasa2017mind}.

Operationally, each time we \emph{notice} distraction and \emph{return} to the target, we add a data point that reinforces its priority (importance) through selection history and automatization \cite{awh2012, logan1988}.
Over many iterations, these micro-updates make the target increasingly default, reducing control demands during sustained work \cite{shiffrin1977, langner2013}.

A complementary lever is \emph{interest}.
Cultivating curiosity toward the object of focus increases its intrinsic reward value and biases attention and memory systems toward it \cite{gottlieb2013, gruber2014, kang2009}.
Contemporary mindfulness protocols explicitly leverage interest/curiosity to reshape reward-based habits that compete with sustained attention \cite{brewer2013}.

\paragraph{Actionable levers for training stable attention}
\begin{itemize}[leftmargin=*]
  \item \textbf{Train importance via repetition.} Use brief, frequent focus--return cycles (e.g., breath or single-target tasks) to strengthen selection history and automatize the on-task state \cite{shiffrin1977, logan1988, awh2012, maclean2010}.
  \item \textbf{Train interest via curiosity.} Adopt an investigative stance (``what is \emph{actually} present here?'') to raise intrinsic value and improve encoding/retention for the current target \cite{gottlieb2013, gruber2014, kang2009, brewer2013}.
\end{itemize}

Together, these levers align with established mechanisms of attention, motivation, and plasticity: repetition shifts control from effortful to automatic, while curiosity increases the target's reward value---jointly stabilizing focus during high-uncertainty work \cite{pessoa2009, dayan2011, langner2013}.

% 18_success_factors.tex — Common Success Factors for Nobel and Turing Awards

\section{Common Success Factors for Both Nobel and Turing Awards}
\label{sec:success_factors}

Analysis of Nobel and Turing laureates reveals a shared set of structural and personal characteristics that are highly correlated with ultimate success.
These factors underscore the idea that while exceptional individual ability is necessary, it is actualized within the highly structured social system of science \citep{zuckerman1977}.

\subsection{Foundational Career Choices (Strategic)}
\begin{itemize}[leftmargin=*]
    \item \textbf{Apprentice Under a Master.} The most significant predictor of future success is the choice of doctoral and postdoctoral advisor.
    Zuckerman's seminal study found that nearly half of all U.S.\ Nobel laureates had trained under a previous laureate, a phenomenon she described as a form of ``apostolic succession'' where tacit knowledge, research ``taste,'' and high standards are transmitted \citep{zuckerman1977}.
    This pattern of elite mentorship is also prevalent among Turing laureates \citep{devanbu2016}.

    \item \textbf{Embed in an Elite Institutional Environment.} A disproportionately high number of laureates are concentrated at a small handful of elite universities (e.g., MIT, Stanford, Berkeley, Harvard, Princeton) and, for computer science, corporate research labs.
    These institutions provide the necessary concentration of talent, resources, and intellectual ferment required for breakthrough research \citep{devanbu2016}.
\end{itemize}

\subsection{Research Strategy and Execution (Skill-Based)}
\begin{itemize}[leftmargin=*]
    \item \textbf{Master the Art of Problem Selection.} Laureates excel at identifying problems that are not merely interesting but fundamental, and which are ``ripe'' for a solution.
    This involves avoiding incrementalism in favor of high-risk, high-reward questions that have the potential to open new fields of inquiry.

    \item \textbf{Adopt a Long-Term Horizon.} Prize-winning research is almost never a short-term project.
    It requires a sustained, often decade-long, focus on a single deep problem.
    This perseverance is a defining characteristic of laureates who resist pressures for constant, short-term publication outputs.
\end{itemize}

\subsection{The System of Cumulative Advantage (The Matthew Effect)}
Both career paths are profoundly shaped by what Robert K.\ Merton termed the ``Matthew Effect,'' where early advantages and recognition are compounded over time \citep{merton1968}.
\begin{itemize}[leftmargin=*]
    \item An early publication in a high-impact journal, an association with a top lab, or a prestigious fellowship provides a scientist with superior resources and visibility.
    This, in turn, increases their chances of future success, creating a self-reinforcing cycle where the ``rich get richer'' in terms of scientific capital \citep{merton1968}.
    For an aspiring researcher, this highlights the disproportionate importance of early-career strategic decisions.
\end{itemize}

\section{Key Differences in Prize-Winning Contributions}
\label{sec:prize_differences}

While the social dynamics are similar, the nature of the recognized work reveals the distinct philosophical foundations of the natural sciences versus computer science.

\subsection{The Nobel Prize Path: A Focus on Discovery}
The Nobel Prizes in Physics, Chemistry, and Medicine are awarded for \textbf{discovering} a fundamental aspect of the natural universe.
The work reveals a pre-existing reality.
\begin{itemize}[leftmargin=*]
    \item \textbf{Contribution Type: Uncovering Reality.} The laureate uncovers a physical law, a biological mechanism, or a chemical structure that was always present but previously unknown to humanity (e.g., the structure of DNA, the existence of gravitational waves).
\end{itemize}

\subsection{The Turing Award Path: A Duality of Discovery and Invention}
The Turing Award honors work that \textbf{creates} new conceptual or functional realities.
It is often about building what \textit{can be}, not just finding what \textit{is}.
\begin{itemize}[leftmargin=*]
    \item \textbf{Contribution Type 1: Inventing a Foundational Artifact.} A significant portion of Turing Awards recognize the design and implementation of entirely new systems that became pillars of modern computing.
    The creation of the UNIX operating system at Bell Labs is a classic example of an artifact whose practical and conceptual impact was deemed of ``lasting and major technical importance'' \citep{shasha2003}.

    \item \textbf{Contribution Type 2: Discovering a Law of Computation.} This theoretical path mirrors the Nobel ethos, where laureates discover or formalize fundamental, abstract truths about computation itself, such as the theory of NP-completeness.
    This work defines the boundaries of what is computationally possible.
\end{itemize}

% 19_tacit_knowledge.tex — The Role of Tacit Knowledge in Skill Transmission

\section{The Role of Tacit Knowledge in Skill Transmission}
\label{sec:tacit_knowledge}

A central finding in the study of elite scientists is that formal education alone is insufficient to produce world-class researchers \citep{zuckerman1977}.
Mentors transmit a vast body of unwritten, experience-based skills known as \textbf{tacit knowledge}.
First articulated by scientist-philosopher Michael Polanyi, tacit knowledge is encapsulated in his famous aphorism, ``we can know more than we can tell'' \citep{polanyi1966}.
It is the intuitive ``operating system'' of an expert practitioner, learned not from textbooks but through direct apprenticeship, observation, and imitation.

This transmission of tacit knowledge is the primary mechanism behind the ``apostolic succession'' observed in Nobel and Turing lineages.
Below is a systematic breakdown of the key components of this knowledge.

\subsection{The Art of Problem Selection (Research ``Taste'')}
This is the ability to identify a problem that is not just solvable, but also important and impactful.
It includes:
\begin{itemize}[leftmargin=*]
    \item \textbf{Developing a ``Nose'' for a Good Problem.} An intuition to distinguish between a dead end, an incremental advance (``salami-slicing''), and a problem that could open a new field.
    This is learned by observing a mentor evaluate and debate research directions \citep{zuckerman1977}.
    \item \textbf{Sensing ``Ripeness''.} Knowing not just \textit{what} to work on, but \textit{when}.
    A problem may be important, but if the necessary tools (e.g., computational power, imaging technology) are not yet available, it is a poor choice.
    This sense of timing is learned by watching a mentor's strategic decisions \citep{medawar1979}.
    \item \textbf{Calibrating Ambition.} Scoping a problem to be significant enough for a major contribution but narrow enough to be achievable within a realistic timeframe.
\end{itemize}

\subsection{Tacit Skills in Experimental Practice \& Execution}
This is the ``craft'' of science---the hands-on knowledge that separates a novice from a master.
\begin{itemize}[leftmargin=*]
    \item \textbf{The ``Feel'' of an Experiment.} In a wet lab, this is the physical intuition for a procedure---the ``look'' of healthy cells or the subtle sound of a properly working vacuum pump.
    In computational work, it's the intuition for when a model is ``behaving weirdly'' before it fully crashes.
    This is learned by doing alongside an expert \citep{polanyi1966}.
    \item \textbf{Troubleshooting Intuition.} When an experiment fails, an experienced scientist has a probabilistic mental map of the most likely culprits.
    An apprentice learns this flowchart by debugging failures \textit{with} their mentor.
    \item \textbf{Knowing the ``Real'' Standard of Evidence.} Understanding what it \textit{actually} takes to convince a skeptical reviewer in a specific subfield, which often goes beyond textbook statistical significance.
    This is absorbed by seeing which results a mentor deems ``clean enough'' to publish \citep{zuckerman1977}.
\end{itemize}

\subsection{Tacit Skills in Interpretation \& Communication}
Data is useless until it is interpreted and communicated effectively as a compelling narrative.
\begin{itemize}[leftmargin=*]
    \item \textbf{Building a Compelling Narrative.} A research paper is a persuasive story, not a chronological log.
    A mentor teaches how to frame results, what to emphasize, and how to anticipate and address potential criticisms.
    This is learned primarily through the process of co-writing \citep{medawar1979}.
    \item \textbf{Navigating Peer Review.} The ability to decode reviewers' comments---distinguishing a polite suggestion from a fatal flaw---and to write a response letter that is both respectful and firm.
    \item \textbf{Presenting with Authority.} This includes reading the audience during a talk, knowing which details to skip if time is short, and gracefully answering challenging questions.
\end{itemize}

% 20_junior_phd_guide.tex — Playing the Long Game: A Junior PhD's Guide

\section{Playing the Long Game: A Junior PhD's Guide to Navigating Academia}
\label{sec:junior_phd_guide}

A few years into a PhD, the initial excitement often gives way to a nagging question: \textit{Am I doing this right?}
The roadmap for a graduating student is clear---job packets, interviews, negotiations.
But for those of us still in the thick of it, the path is less defined.
The advice that follows is a distillation of lessons learned the hard way about navigating the hidden curriculum of academia: the unspoken rules of prestige, the true meaning of metrics, and the art of building a career that lasts.

\subsection{The Reality of a Rigged Game: Branding and Your Niche}

Let's start with a truth that makes many uncomfortable: your institution's brand and your advisor's reputation are immensely powerful.
The data is unequivocal---a small handful of elite universities form a dense, self-reinforcing hiring network, placing their graduates into a disproportionate share of faculty jobs.
It's easy to see this and conclude the game is rigged.

And in a way, it is.
So, what's the move if you're not starting from inside this network?
You don't play their game.
You change the rules.
The counter-strategy to prestige is to become the indispensable expert in a well-defined niche.
Don't be another voice in a crowded field; create your own field.
Find a novel problem, develop a unique dataset, or pioneer a new evaluation method.
When you become the person synonymous with an important question, prestige follows your work, not the other way around.
It's better to be the leading expert in a smaller, defensible area than a minor player in a saturated one.

\subsection{Metrics That Matter: It's About the Trajectory, Not the Total}

Yes, citations and your h-index matter.
Anyone who tells you otherwise is either being naive or disingenuous.
But they don't matter in the way most junior students think.
When senior faculty or hiring committees look at your Google Scholar, they aren't just counting.
They are looking for a story, a trend---the classic ``hockey stick'' curve that shows your impact is accelerating.

A flat trajectory, even at a high level, is less compelling than a steepening curve.
It signals momentum and suggests your work is starting to open new doors for others.
Don't fall into the ``paper machine'' trap, churning out incremental work to boost your publication count.
This strategy often backfires.
A long list of papers without a coherent, impactful narrative can be a red flag.
Instead of asking, ``Will this get published?'' ask a better question: \textbf{``What new question will this allow someone else to ask?''}
That is the work that builds a trajectory.

\subsection{The ``Slam Dunk'' Paper: Quality Over Quantity}

In the publish-or-perish culture of academia, it feels like the goal is to accumulate publications.
But impact is not cumulative; it's often exponential.
A single, high-impact ``slam dunk'' paper can do more for your career than a dozen incremental ones.
It puts you on the map, gives your research a clear anchor, and tells a memorable story.

This requires patience and strategic courage.
It means resisting the urge to ``slice a project thin'' for a quick publication.
Sometimes, it means pushing back and asking for more time to let a promising idea reach its full potential.
The guiding question, borrowed from David Forsyth, is a powerful one: \textbf{``What problem do I need to solve to become famous?''}
While ``famous'' may be hyperbole, the strategic core is sound: What problem, if you solved it, would make an undeniable and lasting impact?
Aim for that.

\paragraph{What Makes a Paper a ``Slam Dunk'' Paper?}

A ``slam dunk'' paper is one that leaves a lasting mark on the field.
Several common properties can be identified:

\textbf{Immortal results.} These are contributions that remain relevant for decades.
They often take the form of benchmarks, new concepts, or methods that fundamentally change how core tasks are approached.

\textbf{Timeliness.} The impact of a paper is amplified when it addresses questions that are especially pressing at the moment of publication, resonating with the ongoing discourse of the research community.

\textbf{Wide interest.} Great papers go beyond a narrow niche and speak to a broad audience.
Their ideas and results are accessible and useful to many researchers across subfields.

Many well-known works exemplify these properties.
Below, we highlight a few such ``slam dunk'' papers and explain why they achieved such status.

\cite{mirage} challenged the prevailing belief that emergent capabilities in LLMs appear only at scale.
Instead, it demonstrated that these effects were artifacts of evaluation metrics rather than fundamental properties.
This work is considered a slam dunk because it overturned a widely accepted assumption with clear, accessible insights and strong empirical evidence.
It reshaped how researchers perceive emergent capabilities in LLMs, reached a broad audience across the AI community, and will remain relevant whenever scaling effects are studied.

\cite{imagenet} is another example.
At the time, computer vision research suffered from a lack of large-scale labeled datasets.
ImageNet filled this gap by providing an unprecedented dataset that enabled a new wave of experiments.
It quickly became the foundational benchmark for vision research, making it indispensable for nearly anyone working in the field.

\cite{batchnorm} introduced the simple but powerful idea of normalization, dramatically improving the speed and stability of neural network training.
Because training stability is central to nearly all of machine learning, this contribution became immortal, shaping the methodology of countless later works.

\cite{resnet} proposed residual learning, enabling the successful training of networks with hundreds of layers.
This insight made deep architectures practical, influenced a wide research community, and remains central in modern architectures.
Its simplicity and broad applicability secured its place as one of the most impactful papers in deep learning.

\subsection{Fear of Being Scooped is a Signal to Change the Game}

The fear of being scooped is universal.
But instead of a source of anxiety, learn to use it as a strategic signal.
If your idea feels easily scoopable, it's because you're in a crowded race where everyone is running on the same track.
You have two choices: be the fastest runner, or find a different track.

If you aren't the fastest coder or the best-resourced lab, don't compete on speed.
Compete on taste, on rigor, on a unique perspective.
Harden your idea.
Raise the bar on your experiments so high that casual competitors won't follow.
Make your work definitive.
Let the fear of being scooped guide you away from the frantic, shallow races and toward the deep, defensible problems where you can build something that lasts.

\subsection{Building Your Support System: Advisors and Community}

Good work is necessary, but it is not sufficient.
Your ideas need advocates, and your career needs a community.
This doesn't happen by accident; you build it deliberately.

\begin{itemize}[leftmargin=*]
    \item \textbf{Manage Your Advisor.} Your advisor is your most important ally, but they are not a mind reader.
    Be direct.
    State your goals explicitly (``My goal is a top academic job...'') and ask for concrete advice (``...what one or two things should I focus on in the next six months?'').
    Schedule dedicated feedback meetings after milestones like a paper submission, when the work is a complete package and the feedback can be focused and actionable.
    \item \textbf{Build Your Champions.} Ask for advice even when you think you have a solid plan.
    Approach mentors and senior colleagues with your strategy and ask them to identify your blind spots.
    This not only strengthens your work but also invites them into your journey, turning them from passive colleagues into invested champions who will advocate for you when you're not in the room.
    \item \textbf{Create the Community You Need.} Service is not a tax on your time; it's an investment in the world you want to work in.
    If you see a gap---a lack of visibility for junior researchers, a need for a specific workshop---fill it.
    Starting a speaker series or a reading group does more than just look good on a CV.
    It builds your network, creates a surface area for your ideas, and establishes you as a leader.
\end{itemize}

% 21_academic_positioning.tex — Positioning for Academic Success Under Structural Constraints

\section{Positioning for Academic Success Under Structural Constraints}
\label{sec:academic_positioning}

While the academic enterprise is often idealized as a pure meritocracy, the reality is shaped by institutional prestige and advisor networks~\citep{bhattad2025jrgrads}.
Hiring data consistently reveal an outsized influence of a small set of institutions and prominent mentors, reinforcing cycles of elite selection~\citep{clauset2015systematic}.
Given this structural inequality, junior researchers outside privileged networks must intentionally craft strategies to establish visibility and influence.
Here we synthesize a systematic approach, grounded in Bhattad's candid guidance~\citep{bhattad2025jrgrads}, focusing on developing \emph{defensible niches}, generating \emph{anchor results}, strategically leveraging mentors, deliberately building academic community, and proactively safeguarding personal well-being.
Each recommendation operationalizes principles directly complementary to our \textbf{Velocity--Vectoring} framework: clarifying high-impact uncertainties, testing critical assumptions through rapid iteration, and continuously updating strategic direction.

\paragraph{Establishing a Defensible Research Niche.}
Prestige-driven academic hiring often privileges institutional affiliation.
To mitigate these constraints, junior researchers must define a compelling niche---articulating a distinct problem space, unique evaluation methodology, or specialized dataset or artifact---that shifts attention toward their work.
Rather than competing on terms favoring elite institutional prestige, researchers gain leverage by strategically redefining the game around their distinctive strengths~\citep{bhattad2025jrgrads}.
Establishing a robust niche positions the researcher as the default point of reference in their chosen specialization, thus redirecting prestige from institutions to individuals and ideas.

\paragraph{Optimizing Impact Trajectory over Absolute Metrics.}
Contrary to conventional wisdom, raw citation counts or paper quantities alone rarely ensure long-term career success.
Instead, influential academic careers are frequently characterized by accelerating impact, or a ``hockey-stick'' trajectory of citations and recognitions~\citep{bhattad2025jrgrads}.
To engineer such growth, researchers must prioritize projects that produce foundational questions, enabling other scholars to build extensively upon their contributions.
The key strategic heuristic becomes not simply ``will this publish?'', but rather ``what new directions will this open for future researchers?''

\paragraph{Crafting Anchor Papers for Exponential Visibility.}
Incremental publications produce linear growth, but a single landmark paper---an ``anchor''---can yield exponential reputational returns~\citep{bhattad2025jrgrads}.
Such work fundamentally shifts conversations and attracts sustained attention.
The pursuit of anchor papers often demands resisting pressures to publish prematurely; instead, scholars should strategically invest additional time to solidify results into enduring contributions.
Identifying potential anchors early, explicitly defining stringent evaluative criteria, and carefully scoping accordingly are critical steps.

\paragraph{Using Competition and Scoops as Strategic Signals.}
The fear of being ``scooped'' can indicate a crowded intellectual marketplace~\citep{bhattad2025jrgrads}.
In response, scholars can either accelerate execution (shortening iterative cycles, tightening experimental loops), or, more strategically, redefine the competitive arena by adopting rigorous, innovative evaluation standards or leveraging unique data assets inaccessible to others.
Choosing battles where competitors lack incentives or resources to follow ensures sustainable differentiation.

\paragraph{Leveraging Advisor Relationships Strategically.}
Advisors are essential resources, yet junior researchers frequently underutilize mentorship relationships by failing to clearly articulate their ambitions.
A proactive approach---explicitly stating long-term goals (e.g., faculty roles) and regularly soliciting targeted, milestone-specific strategic advice---transforms advisors from passive guides into active career advocates~\citep{bhattad2025jrgrads}.
Structured meetings post-submission or post-review cycles provide particularly effective opportunities for high-impact advisor input.

\paragraph{Building Visibility and Community Deliberately.}
Visibility does not spontaneously emerge from excellent work alone; it must be consciously constructed~\citep{bhattad2025jrgrads}.
Effective strategies include frequent practice talks in low-stakes environments, disciplined question-asking at public seminars, thoughtful follow-ups with speakers, and consistent, concise social media dissemination (e.g., structured threads or ``TwePrints'').
Over time, these activities cultivate professional recognition, facilitating collaboration and invitations to higher-visibility opportunities.

\paragraph{Designing Service as Strategic Investment.}
Academic service should transcend obligation and function as community-building aligned to researchers' genuine interests~\citep{bhattad2025jrgrads}.
Creating forums such as specialized workshops, career seminar series, or innovative benchmarking platforms not only fills unmet community needs but also amplifies personal visibility and leadership within a targeted intellectual niche.

\paragraph{Protecting Sanity and Ensuring Sustainability.}
The inherent uncertainties of research produce frequent setbacks, contributing to widespread experiences of imposter syndrome~\citep{bhattad2025jrgrads}.
Deliberately choosing problems suited to individual intellectual strengths and taste, rather than purely prestige-driven or overly crowded fields, can reduce unnecessary stress.
Additionally, proactive strategies---such as regular, intentional rest periods, selective workload management, and maintaining a reflective record of incremental accomplishments (``small wins'')---are critical to sustaining productivity and well-being.

\paragraph{Summary of Actionable Recommendations:}
To concretely implement these principles, junior scholars can adopt the following structured quarterly practices:

\begin{itemize}[leftmargin=*]
  \item Draft a two-page ``niche memo'' defining core problems, distinctive evaluation strategies, and essential artifacts.
  \item Clearly delineate at least one potential anchor paper, explicitly specifying evaluative criteria and essential experiments.
  \item Regularly raise methodological rigor by adding stress tests or robustness checks inaccessible to casual competitors.
  \item Consistently release artifacts explicitly optimized for ease of external adoption (e.g., clearly documented code, reproducible benchmarks).
  \item Hold structured advisor check-ins after key milestones, explicitly requesting focused strategic advice.
  \item Maintain a systematic visibility routine: monthly presentations, thoughtful seminar engagement, and regular digital outreach.
  \item Establish and sustain at least one regular community-building initiative closely aligned with research interests.
  \item Track a simple trajectory log of external adoption signals---citations, artifact usage, speaking invitations---to gauge strategic effectiveness.
  \item Conduct weekly brief sprints identifying and experimentally testing the highest-impact uncertainties.
  \item Enforce clear ``go/kill'' criteria based on external uptake, clarity of contribution, and incremental insight generation.
\end{itemize}

Collectively, these practices integrate into our broader Velocity--Vectoring framework, systematically addressing the structural constraints of academic success, strategically positioning researchers for enduring impact, and proactively fostering both intellectual and personal sustainability.

% 22_dual_strategies.tex — Dual Strategies for AI Research: Insights from MIT Lecture 1.2

\section{Dual Strategies for AI Research: Insights from MIT Lecture 1.2}
\label{sec:dual_strategies}

Paul Liang, MIT. Lecture 1.2 - Introduction to AI Research

Links:
\href{https://mit-mi.github.io/how2ai-course/spring2025/}{Course website}
\href{https://mit-mi.github.io/how2ai-course/spring2025/schedule/lec1.2%20-%20AI%20research.pdf}{Slides}
\href{https://www.youtube.com/watch?v=104FX8MYKAM}{Video}

Overview
\begin{itemize}
    \item How to generate research ideas
    \item How to read a paper
    \item How to execute research ideas
    \item How to write a paper
\end{itemize}

In a foundational lecture from the MIT ``How to AI Almost Anything'' course (Spring 2025), the instructor presents a systematic process for AI research.
This process is framed as an iterative loop of (1) initial observation and ideas, (2) literature review, (3) research questions and hypotheses, (4) test with experiments, and (5) analyze data and (6) report conclusions.
The lecture's core advice for navigating this loop centers on two distinct, risk-aware strategies for generating research ideas.

One of the most important questions in scientific research is how to get research ideas.
There are two general ways, the first one is to turn a concrete understanding of research's failings to a higher-level experimental question, which is called the \textbf{``bottom-up discovery''}.
It is a great tool for incremental progress, but it may preclude breakthroughs that bridge larger gaps.

Another way is to move from a higher-level question to a lower-level concrete testing of that question, which is called the \textbf{``top-down discovery''}.
It favors bigger ideas, but can be disconnected from reality.

\textbf{Research questions.} One or several explicit questions regarding the thing you want to know.
Hypotheses are easier to draft with ``yes-no'' questions rather than ``how to'' questions.

\textbf{Hypothesis.} What you think the answer to the question may be a-priori.
Should be \textit{falsifiable}: if you get a certain result the hypothesis will be validated, otherwise disproved.

Examples of research questions.

Beware ``Does X Make Y Better'' ``Yes''

The above question/hypothesis is natural, but indirect.
If the answer is ``no'' after your experiments, how do you tell what's going wrong?

Usually you have an intuition about \textit{why} X will make Y better (not just random).

Can you think of other research questions/hypotheses that confirm/falsify these assumptions?

How to do Literature Review and Read a Paper?
\begin{enumerate}
    \item Google scholar
    \item Papers with code, Github, Huggingface
    \item Recent conference proceedings
    \item Blog posts
    \item Survey papers, tutorials, courses.
\end{enumerate}

Derisking testing of research ideas:
\begin{enumerate}
    \item Gather and process dataset, visualize data, gather labels, do data splits.
    \item Implement the most simple pipeline and get it working. Pipeline = data loading + basic model + eval function + loss/visualization/deployment.
    \item Change one component of the model at a time, repeat x10 (top-down or bottom-up).
    \item Find what works best, and exploit.
    \item Scale up experiments, repeat across multiple datasets.
    \item Careful ablation studies.
    \item Qualitative comparisons and visualizations.
    \item Repeat until successful.
\end{enumerate}

\paragraph{Bottom-Up Discovery (Low-Risk)}
This strategy is presented as the ``safer'' path, ideal for ``incremental and steady progress.''
Researchers employing this method begin by surveying state-of-the-art (SOTA) models to identify not just what works, but more critically, what fails.
The core of this approach is rigorous \textbf{error analysis} to taxonomize these failure modes.
The subsequent research goal becomes proposing and implementing a novel method to specifically remediate these identified errors.
This bottom-up approach provides a structured timeline for a project: a proposal to scope the area, a midterm report to analyze existing failures, and a final report to present the novel solution.

\paragraph{Top-Down Research (High-Risk)}
In contrast, top-down research begins with ``larger research themes'' or ``moonshot ideas.''
Instead of starting from existing failures, this approach starts from a high-level vision that, if achieved, would represent a ``large leap'' in the field.
The instructor notes this strategy is inherently ``more risky.''
Success requires breaking the ambitious goal into feasible sub-problems and engaging in frequent brainstorming with mentors to validate the direction and ensure the project is tractable within a given timeframe.

\paragraph{De-risking Hypotheses via Falsifiable Questions}
A key tactical recommendation for de-risking \textit{any} research idea lies in the rigorous formulation of the research question.
The instructor strongly advocates for phrasing questions to be \textbf{``yes/no in nature.''}
This structure is preferred because it makes hypotheses ``easily verifiable and falsifiable.''
This is contrasted with less effective questions, such as ``Does X make Y better?''
If the answer is no, the research path is blocked.
A more robust formulation is \textbf{``\textit{Why} X makes Y better?''}
This question tests a specific, falsifiable \textit{property} of X, providing a clear path forward and a deeper insight regardless of the outcome.


% 97_acknowledgments.tex — Acknowledgments

\subsubsection*{Acknowledgements}
This research was funded in part by a Sloan Fellowship, NSF 2046795, 1909577, 1934986, and NIFA award 2020-67021-32799.
This research was funded in part by the Stanford School of Engineering fellowship and EDGE scholar fellowship from Stanford University.
This research was also supported by the Stanford CS summer internships CURIS internship program.


\clearpage
\bibliography{references}
\bibliographystyle{neurips}

\clearpage
\appendix
% 98_appendix.tex — Appendix

\section{Appendix}
\label{appendix:main}

\subsection{Rebuttal Checklist}
\label{appendix:rebuttal_checklist}

A one-page companion to Section~\ref{sec:rebuttals_and_ac_comments}, intended to be run in the hour before an author response is posted.

\paragraph{Order of work.}
\begin{enumerate}
    \item Answer the AC metareview first, in the AC's order and numbering.
    \item Then the lowest-scored review, then the middling ones.
    \item Leave a strong accept largely alone; check its concession density.
    \item Draft the AC comment last, closing with an enumerated list of manuscript changes and their locations.
\end{enumerate}

\paragraph{Substance.}
\begin{itemize}
    \item Does every weakness (W1, W2, \ldots) and question (Q1, Q2, \ldots) have a visible answer?
    \item Is every number findable by the reviewer in a named table, section, or line?
    \item Has any technical premise been conceded without first being verified?
    \item Is any concession worded more strongly than the reviewer's own framing?
    \item Does every concession state what changes and where, and sit beside what survives it?
    \item Is each post hoc comparison labelled as a sensitivity check, with $n$ stated?
    \item Are association and causation kept distinct, with the identifying design named where only association is available?
    \item Would a co-author who knows the paper follow every claim?
    \item Does the response end on what stands rather than on a tally of what changed?
\end{itemize}

\paragraph{Voice.}
\begin{itemize}
    \item Sentence lengths varied, not uniform.
    \item Roughly one em-dash per section.
    \item No ``not $X$, but $Y$,'' ``it is worth noting that,'' ``this is precisely why,'' ``we would like to emphasise,'' ``importantly,'' ``furthermore,'' ``crucially,'' ``delve,'' ``leverage,'' or ``robust'' as filler.
    \item No triads used for rhythm rather than content.
    \item No two consecutive paragraphs opening the same way.
    \item Concrete nouns and specific numbers in place of abstractions.
    \item Read aloud once, in full.
\end{itemize}

\paragraph{Mechanical pass.}
\begin{itemize}
    \item Strip \LaTeX{} escaping and convert to the submission system's markdown.
    \item Confirm every header sits on its own line.
    \item Check the character count against the limit.
    \item Confirm nothing was truncated, especially the final paragraph.
\end{itemize}


\end{document}
